\chapter{Experiment 2: Virtual Wall Constraint}
\label{app:exp2}

\section{Overview}

Experiment~2 evaluates compliant interaction with a virtual plane.
The wall is defined by $\hat{\mathbf{n}}=[0,0,1]^\top$, so the normal
direction coincides with the base-frame $z$-axis and $\mathbf{W}=\mathbf{I}_3$.

\textbf{Why the matrices are diagonal for this specific case:}
Because $\hat{\mathbf{n}}$ is the $z$-axis and $\hat{\mathbf{t}}_1$,
$\hat{\mathbf{t}}_2$ are the $x$- and $y$-axes, the wall frame
$\mathbf{W} = [\hat{\mathbf{t}}_1\;\hat{\mathbf{t}}_2\;\hat{\mathbf{n}}] = \mathbf{I}_3$.
The base-frame stiffness is therefore
$\mathbf{K}_p = \mathbf{I}_3\,\mathbf{K}_{p,\mathrm{local}}\,\mathbf{I}_3 = \mathbf{K}_{p,\mathrm{local}}$,
which is diagonal.
For any tilted wall ($\mathbf{W}\neq\mathbf{I}_3$), the same procedure
produces a full $3\times3$ \abbr{SPD} matrix with non-zero off-diagonals.

\section{Simulated Results}

\begin{figure}[htb]
\centering
\includegraphics[width=\textwidth]{exp2_wall_trials.pdf}
\caption{Virtual wall simulation: three interaction trials.
         The stiff normal direction ($k=1000\,\mathrm{N/m}$) resists
         penetration; the free in-plane direction ($k=50\,\mathrm{N/m}$)
         allows sliding; free tilt ($k_R=5\,\mathrm{Nm/rad}$) allows
         tool rotation in the plane while pitch is locked
         ($k_R=200\,\mathrm{Nm/rad}$).}
\label{fig:exp2_app}
\end{figure}

\section{C++ Implementation}

\begin{lstlisting}[style=cppstyle,
  caption={Experiment 2: Virtual wall controller with CSV logging},
  label={lst:exp2}]
// Experiment 2: Virtual Wall (n_hat = [0,0,1], W = I)
// In-plane: k_free=50 N/m;  Normal: k_stiff=1000 N/m
// Tilt free (t1,t2): k_R_free=5 Nm/rad;  Spin stiff (n): k_R_stiff=200 Nm/rad

// Wall frame (axis-aligned: W = I_3)
// Stiffness in wall frame (diagonal because W=I)
Eigen::Matrix3d Kp_wall, KR_wall;
Kp_wall <<  50,  0,     0,
             0, 50,     0,
             0,  0, 1000;   // [N/m]
KR_wall << 200,   0,  0,
             0, 200,  0,
             0,   0,  5;   // [Nm/rad]

// Rotate to base frame: Kp = W * Kp_wall * W^T
// Here W = I so Kp = Kp_wall (diagonal).
// For a TILTED wall: W != I and Kp becomes FULL SPD matrix.
Eigen::Matrix3d W = Eigen::Matrix3d::Identity(); // change for tilted wall
Eigen::Matrix3d Kp = W * Kp_wall * W.transpose();  // full SPD in general
Eigen::Matrix3d KR = W * KR_wall * W.transpose();

// Damping: D = 2*sqrt(K) per principal direction
// (eigenvalues of Kp: 50, 50, 1000 => D diag: 14.14, 14.14, 63.25)
Eigen::Matrix3d Dp = W *
    Eigen::Vector3d(2*std::sqrt(50), 2*std::sqrt(50), 2*std::sqrt(1000))
    .asDiagonal() * W.transpose();
Eigen::Matrix3d DR = W *
    Eigen::Vector3d(2*std::sqrt(200), 2*std::sqrt(200), 2*std::sqrt(5))
    .asDiagonal() * W.transpose();

// Control law (same as Experiment 1):
Eigen::Vector3d f = Kp * e_p - Dp * pdot;
Eigen::Vector3d m = KR * e_R - DR * omega;
Eigen::Matrix<double,6,1> F; F.head<3>()=f; F.tail<3>()=m;
Eigen::Matrix<double,7,1> tau = J.transpose() * F;
tau += tau_g + tau_c;

// Log wall-specific quantities to CSV
double d_wall = n_hat.dot(p_e - p0) * 1000.0;  // [mm]
double phi_tilt = std::sqrt(e_R(0)*e_R(0)+e_R(1)*e_R(1)) * 180.0/M_PI;
double phi_spin = std::abs(e_R(2)) * 180.0/M_PI;
csv << time << "," << d_wall << "," << phi_tilt << "," << phi_spin;
for (int i=0; i<7; ++i) csv << "," << tau(i);
csv << "\n";
\end{lstlisting}

\section{Python Simulation Code}

The following (from \texttt{thesis\_plots.py}) simulates all three trials:

\begin{lstlisting}[style=cppstyle, language=Python,
  caption={Python simulation for Experiment~2 (virtual wall trials)},
  label={lst:exp2_sim}]
def simulate_sdof(K, D, M, F_fn, dt=0.001, T=4.0):
    t = np.arange(0, T, dt)
    x, xd = np.zeros(len(t)), np.zeros(len(t))
    for i in range(1, len(t)):
        xdd = (F_fn(t[i-1]) - D*xd[i-1] - K*x[i-1]) / M
        xd[i] = xd[i-1] + dt*xdd
        x[i]  = x[i-1]  + dt*xd[i-1]
    return t, x

# Trial A: 10 N normal push (k=1000 N/m, stiff)
F_A = lambda t: 10.0 if 0.5 < t < 1.5 else 0.0
t_A, d_n = simulate_sdof(1000, 2*np.sqrt(1000), 1.0, F_A)

# Trial B: 2 N in-plane slide (k=50 N/m, free)
F_B = lambda t: 2.0 if t > 0.5 else 0.0
t_B, d_t = simulate_sdof(50, 2*np.sqrt(50), 1.0, F_B)

# Trial C: tilt-to-contact (k_tilt=5 Nm/rad free, k_spin=200 stiff)
tau_tilt = lambda t: 0.5 if 0.5 < t < 2.5 else 0.0
t_C, phi_tilt = simulate_sdof(5, 2*np.sqrt(5), 1.0, tau_tilt)
\end{lstlisting}

\section{Parameter Table}

\begin{center}
\renewcommand{\arraystretch}{1.3}
\begin{tabular}{@{}lllll@{}}
\toprule
Direction & \abbr{DOF} & Stiffness & Damping & Physical effect \\
\midrule
$\hat{\mathbf{t}}_1$ ($x$) & Free  & $50\,\mathrm{N/m}$    & $14.1\,\mathrm{Ns/m}$      & Slide freely \\
$\hat{\mathbf{t}}_2$ ($y$) & Free  & $50\,\mathrm{N/m}$    & $14.1\,\mathrm{Ns/m}$      & Slide freely \\
$\hat{\mathbf{n}}$ ($z$)   & Stiff & $1000\,\mathrm{N/m}$  & $63.2\,\mathrm{Ns/m}$      & Stay on wall \\
Tilt ($\hat{\mathbf{t}}_1$) & Free  & $5\,\mathrm{Nm/rad}$   & $4.5\,\mathrm{Nm\,s/rad}$  & Tool tilts to full contact \\
Tilt ($\hat{\mathbf{t}}_2$) & Free  & $5\,\mathrm{Nm/rad}$   & $4.5\,\mathrm{Nm\,s/rad}$  & Tool tilts to full contact \\
Spin ($\hat{\mathbf{n}}$)   & Stiff & $200\,\mathrm{Nm/rad}$  & $28.3\,\mathrm{Nm\,s/rad}$ & Spinning blocked \\
\bottomrule
\end{tabular}
\end{center}


%% =============================================================
%%  APPENDIX C - EXPERIMENT 3: STIFFNESS VARIATION
%% =============================================================
